\chapter{Procedimientos de operación e integración}
\label{app:operacion}

Este apéndice reúne los detalles operativos del sistema descrito en el Capítulo~\ref{ch:implementacion}. Los procedimientos corresponden a las interfaces y configuraciones documentadas; la confirmación de un servicio no acredita la ejecución física de las acciones.

\section{Consulta del código y documentación}
\label{app:codigo}

El repositorio \url{https://github.com/Biorobotica-UNAM/qwen} contiene el código de generación, entrenamiento, inferencia y comparación experimental~\cite{repositorio_qwen_tesis}. Su \href{https://github.com/Biorobotica-UNAM/qwen/blob/3e7fd13b55375d9e776d7a1352673ecf06e45b2a/README.md}{README} explica instalación, configuración y comandos. La Tabla~\ref{tab:mapa_codigo} ofrece enlaces a la revisión pública consultada, identificada por el commit \texttt{3e7fd13b5537}.

\begin{table}[htbp]
\centering
\small
\caption{Ubicación de los componentes en el repositorio de código.}
\label{tab:mapa_codigo}
\begin{tabularx}{\textwidth}{@{}p{0.38\textwidth}X@{}}
\toprule
\textbf{Componente} & \textbf{Ubicación} \\
\midrule
Generación de datos & \href{https://github.com/Biorobotica-UNAM/qwen/tree/3e7fd13b55375d9e776d7a1352673ecf06e45b2a/qwen_gpsr/generation}{\texttt{qwen\_gpsr/\allowbreak generation}} \\
Contratos, catálogos y normalización & \href{https://github.com/Biorobotica-UNAM/qwen/tree/3e7fd13b55375d9e776d7a1352673ecf06e45b2a/qwen_gpsr/domain}{\texttt{qwen\_gpsr/\allowbreak domain}} \\
Entrenamiento QLoRA & \href{https://github.com/Biorobotica-UNAM/qwen/tree/3e7fd13b55375d9e776d7a1352673ecf06e45b2a/qwen_gpsr/training}{\texttt{qwen\_gpsr/\allowbreak training}} \\
Inferencia y servidor HTTP & \href{https://github.com/Biorobotica-UNAM/qwen/tree/3e7fd13b55375d9e776d7a1352673ecf06e45b2a/qwen_gpsr/runtime}{\texttt{qwen\_gpsr/\allowbreak runtime}} \\
Planificador GPT directo & \href{https://github.com/Biorobotica-UNAM/qwen/tree/3e7fd13b55375d9e776d7a1352673ecf06e45b2a/chatgpt_planner}{\texttt{chatgpt\_planner}} \\
Campañas y procesamiento simbólico & \href{https://github.com/Biorobotica-UNAM/qwen/tree/3e7fd13b55375d9e776d7a1352673ecf06e45b2a/experiments}{\texttt{experiments}} \\
Reglas y adaptador versionados & \href{https://github.com/Biorobotica-UNAM/qwen/tree/3e7fd13b55375d9e776d7a1352673ecf06e45b2a/experiments/vendor}{\texttt{experiments/\allowbreak vendor}} \\
Catálogos GPSR & \href{https://github.com/Biorobotica-UNAM/qwen/tree/3e7fd13b55375d9e776d7a1352673ecf06e45b2a/data/CompetitionTemplate}{\texttt{data/\allowbreak CompetitionTemplate}} \\
Informes y registros & \href{https://github.com/Biorobotica-UNAM/qwen/tree/3e7fd13b55375d9e776d7a1352673ecf06e45b2a/reports}{\texttt{reports}} \\
\bottomrule
\end{tabularx}
\end{table}

Los adaptadores y checkpoints de \texttt{models/} son artefactos locales excluidos del control de versiones; su disponibilidad no se deduce del acceso al repositorio. Los modelos base se obtienen de sus fuentes originales. La integración ROS~2 reside en el espacio de trabajo de Justina, mientras que \texttt{experiments/vendor/} conserva el adaptador y las reglas utilizados por el recorrido experimental. Las dependencias, los datos y los pesos necesarios deben identificarse junto con la configuración de cada ejecución.

\section{Comunicación ROS 2 y operación del sistema}

El coordinador envía la orden en \texttt{Parser.\allowbreak Request.\allowbreak command} y
recupera un string JSON en \texttt{Parser.\allowbreak Response.\allowbreak response}.
Una respuesta aceptada contiene \texttt{input} y \texttt{goals};
\texttt{input} utiliza la entrada normalizada si está disponible y,
en caso contrario, la orden enviada. Ante una excepción, el nodo
devuelve \texttt{status=aborted} y la causa.

El cliente comprueba que haya una lista \texttt{goals} en el nivel
superior antes de llamar al conversor. Este solicita la inserción por
lote y después el inicio de planificación. Con el lanzamiento integrado,
las interfaces corresponden a:

\begin{verbatim}
/planning/semantic_parser/semantic_parser
/planning/goals_to_clips/convert
/planning/clips/assert_facts_batch
/planning/clips/start_planning
/planning/clips/results
\end{verbatim}

La conversión utiliza \texttt{clips\_\allowbreak interfaces/\allowbreak srv/\allowbreak String}:
recibe el JSON en \texttt{request.data} y devuelve un mensaje de
estado en \texttt{response.data}. La inserción por lote utiliza el
tipo \texttt{StringArray}, del mismo paquete. El inicio de
planificación utiliza \texttt{std\_\allowbreak srvs/\allowbreak srv/\allowbreak Trigger}.

Al recibir un lote, el nodo encola una tarea que borra su copia del
plan actual, reinicia el entorno mediante \texttt{reset}, inserta
los hechos recibidos y ejecuta el motor. La solicitud de inicio
encola otra tarea que borra la copia del plan, inserta el disparador
\texttt{(start (name action-\allowbreak planning))} y vuelve a ejecutar el motor.
El reinicio del entorno se realiza en la tarea de inserción por lote;
la tarea de inicio no reinicia el entorno y conserva los hechos
cargados para la planificación.

El servicio por lote devuelve \texttt{successful=True} después de
encolar la tarea. El servicio de inicio devuelve
\texttt{success=True} y el mensaje \texttt{Planning started.\allowbreak } en
ese mismo punto del procedimiento. Estas respuestas confirman la
aceptación de las solicitudes, no la terminación de la inserción ni
de la planificación. Si una tarea falla posteriormente, el trabajador
escribe el error en el registro ROS; el callback que ya respondió no
comunica ese fallo mediante su respuesta original.

El adaptador espera hasta tres segundos por la disponibilidad de cada
servicio dependiente y hasta cinco segundos por su respuesta. Rechaza
una lista vacía y devuelve \texttt{OK} cuando ambas solicitudes
reciben respuestas favorables. Estos límites corresponden a la
comunicación de cada operación y no establecen un tiempo máximo
global de generación del plan.

Después de ejecutar el motor, el nodo busca hechos
\texttt{ros-message}. Si existen, solicita su extracción y publica
la salida no vacía mediante \texttt{std\_\allowbreak msgs/\allowbreak msg/\allowbreak String}, con una
instrucción por línea. La obtención del plan se verifica sobre esa
publicación y su contenido, por separado del estado devuelto por los
servicios. El tópico publicado tiene el nombre relativo
\texttt{clips/results}.

\subsection{Recepción del plan y entrega de instrucciones}

El cliente \texttt{qwen\_\allowbreak to\_\allowbreak plan.\allowbreak py} se suscribe al tópico de planes
antes de enviar la orden. Imprime la respuesta semántica, el estado de
conversión y el plan recibido. Puede ejecutarse desde la raíz del
workspace Justina, con el entorno ROS cargado:

{\small
\begin{verbatim}
ros2 launch planning_tools planning.launch.xml
python3 src/Planning/planning_tools/scripts/qwen_to_plan.py \
  "go to the kitchen and find an apple"
\end{verbatim}
}

El lanzamiento y el cliente se mantienen en terminales independientes.
El servidor Qwen debe estar iniciado y accesible mediante
\texttt{QWEN\_\allowbreak SERVER\_\allowbreak URL}. El launch no inicia el servidor, el listener,
el gestor de instrucciones ni el ejecutor físico.
Las reglas se localizan respecto al directorio de trabajo, por lo que
el arranque desde la raíz del workspace forma parte del procedimiento.

El cliente espera, por defecto, diez segundos por servicio o respuesta
y cinco por el plan. El HTTP del launch integrado utiliza cinco segundos,
frente a los quince del launch semántico aislado. Aumentar la espera del
cliente no modifica el timeout HTTP. La suscripción conserva el último
texto recibido, sin identificador que lo vincule con la orden; por ello,
este flujo se utiliza secuencialmente, con una orden a la vez.

La biblioteca C++ \texttt{Planning} ofrece los clientes
\texttt{parse\allowbreak Semantic\allowbreak Command}, \texttt{listen\allowbreak Gpsr\allowbreak Command} y
\texttt{goalsToClips}. Los dos primeros devuelven objetivos y no inician
CLIPS por sí solos; el coordinador solicita la conversión después.
El servicio de escucha obtiene texto de \texttt{/\allowbreak hri/\allowbreak listener/\allowbreak listen}
y lo envía a Qwen, con tiempos de escucha y traducción sucesivos.
Los errores del nodo distinguen \texttt{gpsr\_\allowbreak unavailable}, para fallos
de transporte y timeout, de \texttt{semantic\_\allowbreak error}, para otras
excepciones. Ningún código acredita por sí solo que la orden sea inválida.

El componente \texttt{plan\_\allowbreak manager\_\allowbreak node} recibe el plan y entrega
una instrucción en cada petición. Conserva nombres absolutos históricos;
para conectarlo con las rutas de \texttt{Planning} debe iniciarse aparte:

{\small
\begin{verbatim}
ros2 run plan_manager_node plan_manager_node --ros-args \
  -r /clips/results:=/planning/clips/results \
  -r /request_instruction:=/planning/request_instruction \
  -r /instruction_msg:=/planning/instruction_msg
\end{verbatim}
}

Asignar un namespace no modifica esos nombres absolutos. El gestor
elimina la numeración textual, avanza al entregar cada instrucción y
responde \texttt{No\allowbreak Plan\allowbreak Assigned} cuando ya no tiene un plan activo.
No espera confirmación de que la acción entregada se haya completado.
El cliente de generación termina en el plan y no llama a los controladores.



\section{Inicio del servidor Qwen}

La inferencia puede ejecutarse en el equipo del robot o en una computadora
remota. La entrada \texttt{server.py} delega en
\texttt{qwen\_\allowbreak gpsr.\allowbreak runtime.\allowbreak server}; mantiene residentes el modelo,
el tokenizador y el normalizador. El cliente ROS envía la orden por HTTP,
mientras que la conversión a hechos y CLIPS permanecen en el entorno de
Justina. Las campañas experimentales utilizan su propio proceso CLIPS.

La ruta del adaptador se selecciona explícitamente al iniciar el servicio.
Por ejemplo, desde la raíz del repositorio:

{\small
\begin{verbatim}
python3 server.py --adapter-path models/nl2cd_qwen3b \
  --host 127.0.0.1 --port 8008
\end{verbatim}
}

Para 7B se utiliza \texttt{models/\allowbreak nl2cd\_\allowbreak qwen7b}, que también es la
ruta predeterminada cuando se omite \texttt{-\allowbreak -\allowbreak adapter-\allowbreak path}.
La dirección de escucha y la autenticación se configuran según el
emplazamiento del cliente; \texttt{127.0.0.1} solo admite conexiones
locales. La clave, si se utiliza, se proporciona fuera de los registros
mediante \texttt{QWEN\_API\_KEY} y el encabezado Bearer.
El servidor carga el adaptador indicado, sin comparar checkpoints al arrancar.



\section{Identidad de modelos, datos y recursos}
\label{app:identidad}

Los modelos base son \texttt{Qwen/\allowbreak Qwen2.\allowbreak 5-\allowbreak 3B} y
\texttt{Qwen/\allowbreak Qwen2.\allowbreak 5-\allowbreak 7B}. Los metadatos conservados de inferencia
identifican las siguientes revisiones, respectivamente:

{\scriptsize
\begin{verbatim}
3B: 3aab1f1954e9cc14eb9509a215f9e5ca08227a9b
7B: d149729398750b98c0af14eb82c78cfe92750796
\end{verbatim}
}

Estas revisiones identifican las bases observadas en las campañas
registradas. El campo \texttt{revision} de los adaptadores es nulo y
los scripts no fijan ese argumento; por tanto, la evidencia no acredita
una selección explícita de esas revisiones durante el entrenamiento.
Los estados de \texttt{Trainer}, argumentos y pesos seleccionados se
conservan como productos distintos, según la
Sección~\ref{subsec:seleccion_checkpoint_qwen}.

El corpus compartido está en \texttt{data/\allowbreak datasets/\allowbreak dataset\_\allowbreak gpsr.\allowbreak jsonl}.
La huella SHA-256 del archivo auditado es:

{\scriptsize
\begin{verbatim}
0695f400370b93abbe97fc403a2acaa9f2d6bc8241667788bd466c6ad7c1c406
\end{verbatim}
}

Se utilizó el mismo corpus en ambos entrenamientos; la huella actual
identifica el archivo conservado, pero no corresponde a una medición
realizada durante cada corrida. La auditoría posterior de las
etiquetas y del prefijo comprobado con CLIPS tampoco sustituye un reporte
original de generación que no se conservó.

En GPT se identifican por separado el modelo solicitado, el informado
en la respuesta, la fecha, el prompt y el esquema. La versión nominal
se complementa con huellas y copias de recursos; el prompt vigente no
se atribuye a todas las campañas históricas. El manifiesto de cada campaña identifica el modelo
y los parámetros efectivos utilizados.



\section{Inventario de software}

El inventario conservado del entorno empleado en el trabajo con 7B
incluye Ubuntu 24.04.5 LTS, Python 3.12.3 y ROS 2 Jazzy. La
Tabla~\ref{tab:entorno_registrado_impl} recoge las bibliotecas registradas.
El inventario se distingue de una captura por ejecución y no se extrapola
a cualquier servidor remoto ni al ajuste de 3B sin su registro asociado.

\begin{table}[htbp]
\centering
\small
\caption{Bibliotecas del inventario de software conservado.}
\label{tab:entorno_registrado_impl}
\begin{tabularx}{\textwidth}{@{}Xl@{}}
\toprule
\textbf{Componente} & \textbf{Versión registrada} \\
\midrule
Controlador NVIDIA & 580.178.04 \\
PyTorch / CUDA de su compilación & 2.9.0+cu130 / 13.0 \\
Transformers & 4.51.3 \\
PEFT & 0.18.1 \\
bitsandbytes & 0.49.2 \\
Accelerate & 1.13.0 \\
Datasets & 4.8.4 \\
Tokenizers & 0.21.4 \\
clipspy & 1.0.5 \\
Cliente OpenAI & 2.8.1 \\
jsonschema & 4.25.1 \\
NumPy & 1.26.4 \\
\bottomrule
\end{tabularx}
\end{table}

